\section{Repository}

Questa sezione descrive le modalità di caricamento del progetto sulla repository GitHub e il flusso operativo dalla redazione alla pubblicazione.

\subsection{Struttura degli strumenti}

\subsubsection{Notion}
Notion viene utilizzato come piattaforma di collaborazione e verifica per i contenuti in fase di produzione. I file vengono redatti in Notion e, dopo la verifica, sono pubblicati sulla repository di documentazione su GitHub.

\subsubsection{GitHub}
Il progetto MiroFish Trading Offline utilizza due repository:

\begin{enumerate}
  \item \textbf{Repository Codice}: \url{https://github.com/eghosa02/mirofish_trading_offline}
  
  Contiene il codice sorgente del progetto. I branch devono seguire la nomenclatura: \texttt{feature\#<numero\_issue>}
  
  \item \textbf{Repository Documentazione}: \url{https://github.com/eghosa02/mirofish_trading_offline_doc}
  
  Contiene la documentazione ufficiale del progetto.
\end{enumerate}

\subsubsection{Nomenclatura dei branch}
Il sistema di versionamento prevede branch con la seguente struttura:
\begin{itemize}
  \item \texttt{feature\#<numero>}: per lo sviluppo di nuove funzionalità
  \item Esempio: \texttt{feature\#42} (per la issue numero 42)
\end{itemize}

\subsection{Procedura operativa}

\subsubsection{Setup iniziale - Clonare le repository}

Per iniziare a lavorare al progetto è necessario clonare entrambe le repository. Eseguire i seguenti comandi:

\subsubsection{Clone della repository di codice}
\begin{verbatim}
git clone https://github.com/eghosa02/mirofish_trading_offline.git
cd mirofish_trading_offline
\end{verbatim}

\subsubsection{Clone della repository di documentazione}
\begin{verbatim}
git clone https://github.com/eghosa02/mirofish_trading_offline_doc.git
cd mirofish_trading_offline_doc
\end{verbatim}

\subsubsection{Creazione di una nuova feature}

\paragraph{Step 1: Identificare il numero della issue}
Ogni nuova feature deve essere collegata a una issue su GitHub. Annotare il numero della issue (es. \#42).

\paragraph{Step 2: Creare il branch locale}
\begin{verbatim}
git checkout main
git pull origin main
git checkout -b feature#<numero_issue>
\end{verbatim}

Esempio per la issue \#42:
\begin{verbatim}
git checkout -b feature#42
\end{verbatim}

\paragraph{Step 3: Sviluppare la feature}
Effettuare le modifiche necessarie sul branch creato. È possibile effettuare commit locali durante lo sviluppo:

\begin{verbatim}
git add .
git commit -m "Descrizione breve delle modifiche"
\end{verbatim}

\subsection{Caricamento dei cambiamenti}

\paragraph{Step 1: Verificare i cambiamenti}
Prima di eseguire il push, verificare lo stato dei file modificati:

\begin{verbatim}
git status
\end{verbatim}

\paragraph{Step 2: Aggiungere i file modificati}
\begin{verbatim}
git add .
\end{verbatim}

Se si desidera aggiungere solo specifici file:
\begin{verbatim}
git add <percorso/file>
\end{verbatim}

\paragraph{Step 3: Eseguire il commit}
\begin{verbatim}
git commit -m "Messaggio descrittivo del commit"
\end{verbatim}

\paragraph{Step 4: Eseguire il push}
\begin{verbatim}
git push origin feature#<numero_issue>
\end{verbatim}

\subsection{Creazione della Pull Request}

\paragraph{Step 1: Accedere a GitHub}
Navigare al repository su \url{https://github.com/eghosa02/mirofish_trading_offline}

\paragraph{Step 2: Creare la Pull Request}
Dopo aver eseguito il push, GitHub mostrerà un pulsante \textbf{Compare \& pull request}. Cliccare su di esso.

In alternativa:
\begin{enumerate}
  \item Andare su \textbf{Pull requests}
  \item Cliccare su \textbf{New pull request}
  \item Selezionare il branch \texttt{feature\#<numero\_issue>} da confrontare con \texttt{main}
  \item Cliccare su \textbf{Create pull request}
\end{enumerate}

\paragraph{Step 3: Compilare i dettagli della Pull Request}
\begin{itemize}
  \item \textbf{Titolo}: Deve essere descrittivo e conciso
  \item \textbf{Descrizione}: Includere:
  \begin{itemize}
    \item Riferimento alla issue (\#<numero>)
    \item Descrizione dei cambiamenti effettuati
    \item Eventuali note importanti
  \end{itemize}
\end{itemize}

Esempio di descrizione:
\begin{verbatim}
Closes #42

## Descrizione
Implementazione della nuova funzione di autenticazione.

## Cambiamenti
- Aggiunto modulo di login
- Aggiunta validazione credenziali
- Aggiunto test unitari

## Note
Testato su Windows 10
\end{verbatim}

\paragraph{Step 4: Sottomissione e review}
\begin{enumerate}
  \item Cliccare su \textbf{Create pull request}
  \item Attendere la review del team
  \item Effettuare i cambiamenti richiesti se necessario
  \item Una volta approvata, cliccare su \textbf{Merge pull request}
\end{enumerate}

\subsection{Sincronizzazione e aggiornamenti}

Per mantenere il branch locale aggiornato con i cambiamenti del main:

\begin{verbatim}
git checkout main
git pull origin main
git checkout feature#<numero_issue>
git merge main
\end{verbatim}

Per risolvere conflitti (se presenti):
\begin{enumerate}
  \item Aprire i file con conflitti nell'editor
  \item Risolvere manualmente i conflitti
  \item Effettuare il commit dei cambiamenti
  \item Eseguire il push
\end{enumerate}

\subsection{Best Practices}

\subsubsection{Comandi git essenziali}

Di seguito i comandi principali utilizzati nel flusso di lavoro:

\begin{itemize}
  \item \texttt{git clone <url>} - Clonare una repository
  \item \texttt{git checkout <branch>} - Cambiare branch
  \item \texttt{git checkout -b <branch>} - Creare e cambiare a un nuovo branch
  \item \texttt{git add .} - Aggiungere tutti i file modificati
  \item \texttt{git add <file>} - Aggiungere un file specifico
  \item \texttt{git commit -m "messaggio"} - Effettuare un commit
  \item \texttt{git push origin <branch>} - Inviare i cambiamenti al repository remoto
  \item \texttt{git pull origin <branch>} - Scaricare gli ultimi cambiamenti da remoto
  \item \texttt{git status} - Visualizzare lo stato dei file
  \item \texttt{git log} - Visualizzare la cronologia dei commit
  \item \texttt{git merge <branch>} - Fondere un branch nel branch corrente
\end{itemize}

\subsubsection{Linee guida generali}

\begin{enumerate}
  \item \textbf{Nomenclatura branch}: Usare sempre \texttt{feature\#<numero>} per i nuovi branch
  \item \textbf{Messaggi commit}: Essere descrittivi e concisi
  \item \texttt{Commit frequenti}: Effettuare commit piccoli e logici
  \item \textbf{Pull request}: Creare sempre una PR prima di fare il merge su main
  \item \textbf{Review}: Attendere l'approvazione del team prima del merge
  \item \textbf{Sincronizzazione}: Mantenere il branch sincronizzato con main
  \item \textbf{Notion}: Aggiornare il status della issue su Notion dopo la pubblicazione
\end{enumerate}

\subsubsection{Flusso dalla redazione alla pubblicazione}

\begin{enumerate}
  \item \textbf{Redazione}: Il contenuto viene redatto in Notion
  \item \textbf{Verifica}: Il contenuto è sottoposto a review nel team
  \item \textbf{Approvazione}: Una volta approvato, il contenuto è marcato come pronto per la pubblicazione
  \item \textbf{Pubblicazione}: Il contenuto approvato è pubblicato sulla repository di documentazione via pull request
  \item \textbf{Closure}: L'issue correlata è chiusa dopo il merge
\end{enumerate}
